%==============================================================================
\section{Getting Started --- Bắt đầu với học máy}
%==============================================================================

\begin{dongco}[title={Mục tiêu}]
Phần này mô tả học máy là gì, các loại học máy, một số thuật toán tiêu biểu, use case, và điều kiện tiên quyết để bắt đầu.
\end{dongco}

\subsection{Điều kiện tiên quyết để bắt đầu}

\begin{dongco}[title={Học ML dễ hơn khi chuẩn bị đúng}]
Để bắt đầu học ML, bạn nên nắm một số nền tảng về khoa học máy tính, đồng thời quen với lập trình, thư viện và toán--thống kê.
\end{dongco}

\subsubsection{Ngôn ngữ lập trình: Python hoặc R}

\begin{phuongphap}[title={Gợi ý}]
Có nhiều ngôn ngữ dùng được (C++, Java, Python, R, Julia, \ldots), nhưng Python rất phổ biến trong ML và Data Science.
\end{phuongphap}

\begin{phuongphap}[title={Các chủ đề Python cơ bản nên nắm}]
\begin{itemize}
  \item Variables, basic data types.
  \item Data structures (list, set, dictionaries).
  \item Loops and conditionals.
  \item Functions.
  \item String formatting.
  \item Classes and objects.
\end{itemize}
\end{phuongphap}

\subsubsection{Thư viện và packages}

\begin{phuongphap}[title={Thư viện gợi ý}]
\begin{itemize}
  \item NumPy, Pandas (tiền xử lý).
  \item scikit-learn.
  \item Matplotlib.
\end{itemize}
\end{phuongphap}

\subsubsection{Toán và thống kê}

\begin{phuongphap}[title={Mức toán cần thiết}]
\begin{itemize}
  \item Đại số tuyến tính.
  \item Giải tích cơ bản.
  \item Xác suất--thống kê.
\end{itemize}
\end{phuongphap}

